% ============================================================================
%  sections/methodology/3_3_communication_system_design.tex
%  Section 3.3 -- Communication System Design
% ============================================================================
\section{Communication System Design}

To overcome the coverage limitations of traditional LoRaWAN (single-hop) and the
high control overhead of reactive protocols like AODV, the network adopts the
\textbf{Layered Dynamic Synchronization Energy-saving (LDSE)} protocol.\cite{li2025ldse}
In standard single-hop LoRaWAN, a node positioned several kilometres deep in a
jungle must scale its Spreading Factor up to SF12 to overcome heavy canopy
attenuation, spiking the Time-on-Air and draining the battery with every
transmission. LDSE solves this by structuring the forest into logical, concentric
tiers wrapped around one or more central gateway proxies, so each node
communicates only with its nearest parent tier using short, low-power hops at
optimised spreading factors such as SF7 or SF8.


% ----------------------------------------------------------------------------
\subsection{LoRa Communication Fundamentals}
% ----------------------------------------------------------------------------

LoRa (Long Range) is a proprietary, physical-layer wireless technology developed
by Semtech that enables long-range, low-power, and low-data-rate connectivity for
battery-powered IoT devices. It is based on Chirp Spread Spectrum (CSS)
modulation, a technique in which the carrier frequency is swept linearly over time
to encode information. CSS provides high receiver sensitivity, strong resistance to
interference, and reliable communication under low signal-to-noise ratio (SNR)
conditions, making it well suited to large-scale environmental monitoring
applications such as forest sensing, where nodes are widely dispersed and often
obstructed by dense vegetation.


% ----------------------------------------------------------------------------
\subsection{LoRa Communication Parameters}
% ----------------------------------------------------------------------------

The communication performance of a LoRa link is governed primarily by four
parameters, each involving a trade-off between range, data rate, latency, and
energy consumption.

\begin{itemize}
  \item \textbf{Spreading Factor (SF):} Determines how much a signal is
        ``spread'' in the frequency domain before transmission. A higher SF
        (e.g., SF12) increases communication range and link robustness at the
        cost of a longer packet Time-on-Air and higher energy consumption,
        whereas a lower SF (e.g., SF7) offers faster transmission and lower
        energy use at the cost of reduced range.

  \item \textbf{Bandwidth (BW):} The width of the frequency channel used for
        transmission (commonly 125~kHz, 250~kHz, or 500~kHz in the sub-GHz
        ISM bands). A wider bandwidth increases data rate and reduces
        Time-on-Air but lowers receiver sensitivity, shortening effective range.

  \item \textbf{Coding Rate (CR):} The ratio of forward error correction (FEC)
        redundancy added to a packet. A higher coding rate improves resilience
        to burst interference and multipath fading at the expense of additional
        transmission overhead and airtime.

  \item \textbf{Transmission Power:} The radiated output power of the
        transceiver. Higher transmission power extends range but increases
        energy consumption per packet and, in ISM bands, is often constrained
        by regional duty-cycle and power regulations.
\end{itemize}

Selecting appropriate values for these parameters therefore represents an
important design trade-off between coverage, reliability, latency, and battery
life, and directly motivates the multi-hop, low-SF approach adopted by LDSE
(Section~\ref{subsec:ldse-overview}).


% ----------------------------------------------------------------------------
\subsection{Conventional LoRaWAN Architecture}
% ----------------------------------------------------------------------------

LoRaWAN defines the networking protocol operating above the LoRa physical
layer. A conventional LoRaWAN deployment follows a \textit{Star-of-Stars}
topology, in which multiple end devices communicate directly with one or more
gateways. These gateways transparently forward received packets to a central
Network Server, which performs device authentication, duplicate packet removal,
adaptive data rate management, and routing toward the Application Server.

This architecture offers simple deployment, centralised management, and low
maintenance overhead, making it well suited to applications such as smart
agriculture, smart cities, and general-purpose environmental sensing. However,
communication remains inherently single-hop: every sensor node must maintain a
direct wireless link to a gateway, regardless of terrain or vegetation density.


% ----------------------------------------------------------------------------
\subsection{Limitations of Conventional LoRaWAN}
% ----------------------------------------------------------------------------

Although conventional LoRaWAN provides excellent communication range under
open-terrain conditions, dense forests introduce severe propagation losses due to
vegetation, tree trunks, canopy shadowing, and irregular topography. Maintaining
a direct gateway link under these conditions typically requires raising the
spreading factor, which significantly increases packet airtime, transmission
energy, and channel occupancy. Furthermore, deploying additional gateways to
compensate for coverage gaps substantially increases infrastructure cost --- a
burden that is particularly impractical for low-budget community forest
deployments in Nepal.

These constraints define the coverage-versus-cost trade-off that motivates the
shift away from a single-hop star topology toward an energy-efficient multi-hop
architecture.


% ----------------------------------------------------------------------------
\subsection{LDSE Protocol Overview}
\label{subsec:ldse-overview}
% ----------------------------------------------------------------------------

To overcome the limitations described above, this project adopts the Layered
Dynamic Synchronization Energy-saving (LDSE) protocol. Rather than relying on
long-distance single-hop communication, LDSE establishes an energy-efficient
multi-hop hierarchy in which nearby nodes relay data through synchronised
intermediate parents. This approach extends network coverage, improves link
reliability under forest propagation conditions, reduces transmission energy, and
enables scalable deployment over large monitoring areas while maintaining
low-power operation.

LDSE organises communication through three clear structural phases:
\textbf{Layer Assignment}, \textbf{Dynamic Time-Synchronisation}, and
\textbf{Time-Slotted Contention-Free Data Transmission}, detailed in Sections
\ref{subsec:ldse-architecture}--\ref{subsec:ldse-multichannel}.


% ----------------------------------------------------------------------------
\subsection{Architecture of LDSE}
\label{subsec:ldse-architecture}
% ----------------------------------------------------------------------------

\begin{figure}[htbp]
    \centering
    \includegraphics[width=0.47\linewidth]{ldse_protocol_network_hierarchy.pdf}
    \caption{LDSE protocol network hierarchy}
    \label{fig:ldse-hierarchy}
\end{figure}
\newpage

\textbf{Layered Topology:} The gateway (Layer~0) broadcasts an initialisation
beacon during network setup. Each node assigns itself a hierarchical layer by
incrementing the layer value carried in the beacon it receives:
\begin{equation}
\label{eq:layer_assignment}
L_{\text{node}} = L_{\text{received}} + 1
\end{equation}
This cascading propagation forms a tree-structured topology across the entire
deployment area without any centralised configuration, as shown in
Figure~\ref{fig:ldse-hierarchy}.

When a forest node is more than 2~km from the gateway, alert packets are
automatically forwarded through intermediate relay nodes operating under LDSE
routing, extending effective network coverage without the need for additional
gateways.


% ----------------------------------------------------------------------------
\subsection{Time Synchronization and Epoch Structure}
\label{subsec:ldse-sync}
% ----------------------------------------------------------------------------

\textbf{Time-Slotted Epoch Structure:} LDSE divides time into global operational
epochs, each subdivided into three non-overlapping windows: a
\emph{Synchronisation Window} in which all transceivers wake simultaneously to
receive time-alignment correction beacons from parent layers; a \emph{Data
Transmission Window} implementing TDMA-gated uplinks layer by layer; and a
\emph{Deep-Sleep Period} during which all microcontrollers and radio
transceivers enter ultra-low-leakage power states consuming as little as
$7\,\mu\text{A}$.

\begin{figure}[htbp]
    \centering
    \includegraphics[width=1\linewidth]{ldse_protocol_epoch_structure.png}
    \caption{LDSE Protocol Epoch Structure}
    \label{fig:ldse-epoch}
\end{figure}

\textbf{Time Synchronisation:} The Flooding Time Synchronisation Protocol
(FTSP) distributes a global timestamp from the gateway to all nodes, correcting
clock drift caused by ambient temperature variation in the forest canopy and
ensuring nodes do not miss their relay windows while remaining in deep sleep for
the majority of each cycle.


% ----------------------------------------------------------------------------
\subsection{Parent Selection and Routing Mechanism}
\label{subsec:ldse-parent-selection}
% ----------------------------------------------------------------------------

\textbf{Parent Selection and Implicit Route Exploration (IRE):} When a node
hears multiple candidate parents within a higher tier, it selects the optimal
upstream relay using a combined score that balances link quality against
remaining relay energy:
\begin{equation}
\label{eq:parent_score}
\text{Score} = \alpha \cdot \text{RSSI}_{\text{link}}
+ \beta \cdot \left(1 - \frac{E_{\text{consumed}}}{E_{\text{initial}}}\right)
\end{equation}

Routes are discovered through active beaconing or passive eavesdropping on
existing transmissions, limiting route-discovery overhead to approximately 5\%
of total traffic and avoiding the broadcast storms characteristic of AODV.


% ----------------------------------------------------------------------------
\subsection{Multi-Channel Communication}
\label{subsec:ldse-multichannel}
% ----------------------------------------------------------------------------

\textbf{Multi-Channel Separation:} A public channel carries handshake and
synchronisation messages, while a private channel carries payload data, avoiding
intra-network collisions and congestion at relay nodes.


% ----------------------------------------------------------------------------
\subsection{Power Consumption Model}
% ----------------------------------------------------------------------------

The energy savings of LDSE arise from exchanging extended receiver listening
times for drastically reduced transmission times. The total energy consumed by a
node during one active operational cycle is:
\begin{equation}
\label{eq:total_energy}
E_{\text{total}} = E_{\text{sleep}} + E_{\text{sense}}
+ E_{\text{processing}} + E_{\text{sync}}
+ E_{\text{tx}} + E_{\text{rx}}
\end{equation}
The dominant saving comes from the transmission term. Comparing a standard
single-hop link at SF12 against an LDSE multi-hop link at SF7 over a 3~km
forested path:
\[
  \text{ToA}_{\text{SF12}} \approx 983\,\text{ms}, \qquad
  \text{ToA}_{\text{SF7}}  \approx 56\,\text{ms}
\]
\[
  E_{\text{tx,single}} = 3.3\,\text{V} \times 120\,\text{mA}
                          \times 0.983\,\text{s} \approx 389\,\text{mJ}
\]
\[
  E_{\text{tx,LDSE}}   = 3.3\,\text{V} \times 45\,\text{mA}
                          \times 0.056\,\text{s} \approx 8.3\,\text{mJ}
\]
This represents a transmission energy reduction of approximately $46.8\times$.
Even after accounting for the receiver overhead incurred by relay nodes
($\approx 1.65\,\text{mJ}$ per monitored child slot), the total network energy
savings reach up to \textbf{42\%} compared to classic single-hop topologies,
allowing field-deployed nodes to operate reliably for several years on standard
lithium iron phosphate (LiFePO$_4$) cells.


% ----------------------------------------------------------------------------
\subsection{Forest Radio Propagation Considerations}
% ----------------------------------------------------------------------------

Wireless communication in forest environments is strongly influenced by
vegetation density, foliage absorption, terrain irregularities, and multipath
propagation. Unlike free-space environments, dense forests introduce additional
attenuation that increases with communication distance and canopy density.
Consequently, signal strength degrades more rapidly, reducing communication
reliability for long single-hop links.

During network simulation in ns-3, appropriate propagation models will be used
to represent these environmental effects. The adoption of LDSE mitigates such
attenuation by dividing long communication paths into multiple shorter relay
links, thereby improving packet delivery reliability while reducing the
transmission power and airtime required by individual sensor nodes.


% ----------------------------------------------------------------------------
\subsection{Protocol-Centric Implementation}
% ----------------------------------------------------------------------------

This section outlines how to build and implement a functional two-tier LDSE
protocol stack using an ESP32-S3 host microcontroller paired with an SX1262
LoRa transceiver module.

\begin{figure}[htbp]
    \centering
    \includegraphics[width=0.85\linewidth]{sequential_sync_operational_flowchart.pdf}
    \caption{Sequential Synchronization and Operational Flowchart}
    \label{fig:ldse-flowchart}
\end{figure}


% ----------------------------------------------------------------------------
\subsection{Operational Challenges and Mitigation Strategies}
% ----------------------------------------------------------------------------

Deploying LDSE in a dense forest environment introduces several engineering
challenges, each addressed by a targeted mitigation:

\begin{itemize}
  \item \textbf{Clock Drift due to Thermal Variation:} Diurnal temperature
        swings cause low-cost quartz oscillators to drift by several milliseconds
        over a few hours, risking overlapping TDMA slots and packet collisions.
        This is addressed by inserting a Guard Time Interval
        ($\tau_g = 20\text{--}50\,\text{ms}$) at the edges of each slot and by
        using Temperature-Compensated Crystal Oscillators (TCXO) on custom
        hardware boards to maintain timing alignment across temperature shifts.

  \item \textbf{Relay Congestion during Mass-Alert Events:} During a
        fast-spreading wildfire, many outer nodes may send alerts simultaneously,
        saturating the memory buffers of Layer~1 relay nodes. An asynchronous
        priority queue combined with a dynamic spreading-factor scale-up mechanism
        addresses this: when a relay's buffer exceeds 80\% capacity, it sets a
        congestion flag in its synchronisation beacon, signalling child nodes to
        switch to a higher spreading factor (e.g., SF10) and route alerts
        directly to the gateway, temporarily bypassing the busy relay tier.

  \item \textbf{Hidden Terminal Collisions:} Two sibling nodes in the same
        layer may be out of range of each other yet share the same parent relay.
        If both transmit simultaneously, their signals collide at the relay,
        corrupting both packets. A randomised exponential back-off delay within
        the assigned TDMA slot window, combined with LoRa Channel Activity
        Detection (CAD), which mitigates this: before transmitting, each node listens
        for two symbol durations to detect any active preamble chirp on the
        channel and defers if the channel is busy.
\end{itemize}